% =============================================================================
% 1. INTRODUCCION  (presentacion y estructura: 15 % junto con el aspecto general)
% Espacio: ~0,6 paginas (incluido el esquema).
% =============================================================================
\section{Introducción}
\label{sec:introduccion}

\pendiente{objetivo del trabajo en 2--3 frases: atacar GGH, LWE y Ring-LWE, basados en problemas
  de celosías (SVP, CVP, LWE) que se consideran difíciles incluso para ordenadores cuánticos.}

\pendiente{tareas elegidas y criterio de elección (Tabla~\ref{tab:tareas}). Entorno: Python~3,
  sin librerías de celosías; todo implementado por nosotros con aritmética exacta
  (\texttt{Fraction}). Esquema del trabajo en la Figura~\ref{fig:mapa}.}

\begin{figure}[h]
  \centering
  % Esquema del trabajo: de los datos de cada criptosistema a la solucion, y el nucleo comun.
% Cada fila es un criptosistema. Ajustar los textos cuando esten decididos los metodos
% definitivos de cada tier.
\begin{tikzpicture}[
    paso/.style={draw,line width=0.4pt,text width=2.95cm,minimum height=0.9cm,
                 align=center,font=\footnotesize,inner sep=2pt},
    fila/.style={font=\footnotesize\bfseries,text width=1.3cm,align=left,inner sep=0pt},
    flecha/.style={-{Stealth[length=3.5pt]},line width=0.4pt},
    node distance=0.3cm and 0.45cm]

  \node[fila] (Lg) {GGH};
  \node[paso,right=0.2cm of Lg] (g1) {Base pública $B$\\y objetivo $c$};
  \node[paso,right=of g1] (g2) {LLL:\\base reducida};
  \node[paso,right=of g2] (g3) {Babai /\\enumeración};
  \node[paso,right=of g3] (g4) {Punto $v$\\(SVP / CVP)};

  \node[fila,below=0.95cm of Lg] (Ll) {LWE};
  \node[paso,right=0.2cm of Ll] (l1) {Sistema $(A, b)$,\\$b = As + e$};
  \node[paso,right=of l1] (l2) {Celosía primal\\(embedding)};
  \node[paso,right=of l2] (l3) {LLL: aparece\\el error $e$};
  \node[paso,right=of l3] (l4) {Secreto $s$\\y bit descifrado};

  \node[fila,below=0.95cm of Ll] (Lr) {Ring-LWE};
  \node[paso,right=0.2cm of Lr] (r1) {$(a, b)$ y\\criptograma $(c_1, c_2)$};
  \node[paso,right=of r1] (r2) {Sistema lineal en $\mathbb{Z}$\\(sin reducir mod $q$)};
  \node[paso,right=of r2] (r3) {Resolver y\\redondear: $s$};
  \node[paso,right=of r3] (r4) {$d = c_2 - s\,c_1$:\\palabra de bits};

  \foreach \f in {g,l,r}{
    \draw[flecha] (\f1) -- (\f2);
    \draw[flecha] (\f2) -- (\f3);
    \draw[flecha] (\f3) -- (\f4);}

  % Nucleo comun: recuadro discontinuo bajo las tres filas, del ancho de r1 a r4
  \draw[dashed,line width=0.4pt] ([yshift=-0.4cm]r1.south west) rectangle
                                 ([yshift=-1.1cm]r4.south east);
  \node[font=\footnotesize] at ($([yshift=-0.75cm]r1.south west)!0.5!([yshift=-0.75cm]r4.south east)$)
    {\textit{Núcleo común (propio, exacto):}\enspace
     matrices, Gram-Schmidt, LLL, Babai, enumeración};
\end{tikzpicture}

  \caption{Esquema del trabajo: cada fila es un criptosistema, de los datos de la tarea a la
    solución. Todos los ataques se apoyan en el mismo núcleo de algoritmos de celosías.}
  \label{fig:mapa}
\end{figure}
